\section{Predecessor containment and extremal counterexamples}
\label{sec:paper-predecessor-cloning}

In this section we study a surgery that preserves positive gaps by
copying outgoing neighborhoods under incoming-neighborhood containment.
The surgery itself is valid in an arbitrary finite oriented graph.
Its structural consequences will explicitly state whether they require
minimum order or a stronger lexicographic extremal choice.

We call $D$ \emph{all-deficient} if $g_D(v)>0$ for every vertex.
Such a graph is a counterexample to the second neighborhood conjecture.
The gap-one predecessor and gap-one cycle phenomena occur in the
minimal-counterexample analysis of Brantner et al.~\cite{BrantnerEtAl2009}.
Minimum-order reductions also play a central role in
\cite{EspunyDiazEtAl2025}. The arguments below give statements about
entire incoming classes and establish the precise extremal hypotheses
needed for an independent-module normal form.

\subsection{Copying one row and copying all rows}

All containment predicates and all copied neighborhoods in the next
two results are evaluated in the original graph $D$.

\begin{lemma}[One-row copying under predecessor containment]
\label{lem:paper-single-row-copy}
Let $a\ne b$ and suppose that
$N_D^-(a)\subseteq N_D^-(b)$. Replace the outgoing neighborhood of
$a$ by $N_D^+(b)$, leaving every other outgoing neighborhood unchanged,
and denote the resulting graph by $D^{a\leftarrow b}$. Then this graph
is oriented and
\[
 N_{D^{a\leftarrow b}}^+(a)=N_D^+(b),\qquad
 N_{D^{a\leftarrow b}}^{++}(a)=N_D^{++}(b).
\]
For every $v\ne a$,
\[
 N_{D^{a\leftarrow b}}^+(v)=N_D^+(v),\qquad
 N_{D^{a\leftarrow b}}^{++}(v)\subseteq N_D^{++}(v).
\]
In particular, the new gap at $a$ is $g_D(b)$, all other gaps do not
decrease, and an all-deficient graph remains all-deficient.
\end{lemma}
\begin{proof}
There is no original arc $b\to a$: it would put $b$ in $N_D^-(a)$
and hence in $N_D^-(b)$, contrary to the absence of loops. Nor is there
an original two-step path $b\to x\to a$, since $x\to a$ would imply
$x\to b$ and form a digon with $b\to x$. Thus $a$ is neither a first
nor a second target of $b$.

Every $x\in N_D^+(b)$ is distinct from $a$ and cannot point to $a$,
again because $x\to a$ would imply $x\to b$. The new row at $a$
therefore creates neither a loop nor a digon. The intermediate vertices
in $N_D^+(b)$ retain their outgoing neighborhoods. Their union of
outgoing targets contains neither $a$, by the two-step exclusion just
proved, nor $b$, by the absence of digons. Subtracting the copied first
set $N_D^+(b)$ consequently gives exactly $N_D^{++}(b)$ as the new
second neighborhood of $a$.

For $v\ne a$, only a two-step path using the changed row could produce
a new second target. Such a path has the form $v\to a\to x$, where
$v\to a$ and $b\to x$ are original arcs. Incoming containment gives
the original arc $v\to b$, so $v\to b\to x$ was already an original
two-step path. The first set at $v$ is unchanged, so no new exact
second target appears. Some second targets may disappear. The gap
statements now follow by counting the first and second sets.
\end{proof}

\begin{theorem}[A simultaneous predecessor map]
\label{thm:clone-map}
Let $\rho:V(D)\to V(D)$ satisfy
\begin{equation}
 N_D^-(v)\subseteq N_D^-(\rho(v))\qquad(v\in V(D)).
 \label{eq:paper-clone-containment}
\end{equation}
Define $D^\rho$ by simultaneously setting
\[
 N_{D^\rho}^+(v)=N_D^+(\rho(v))\qquad(v\in V(D)).
\]
Then $D^\rho$ is oriented, and for every $v$,
\begin{equation}
 N_{D^\rho}^{++}(v)\subseteq N_D^{++}(\rho(v)),
 \qquad g_{D^\rho}(v)\geq g_D(\rho(v)).
 \label{eq:paper-clone-map-gaps}
\end{equation}
The map $\rho$ need not be injective, and its image may contain vertices
whose rows are also changed. In particular, all-deficiency is preserved.
\end{theorem}
\begin{proof}
A new loop at $v$ would mean an original arc $\rho(v)\to v$.
By \eqref{eq:paper-clone-containment}, this would imply an original
loop at $\rho(v)$. If $v\to u$ and $u\to v$ were both new arcs,
the original arcs $\rho(v)\to u$ and $\rho(u)\to v$, together with
containment at $u$ and $v$, would give the original label arcs
$\rho(v)\to\rho(u)$ and $\rho(u)\to\rho(v)$. These form a digon
when the two images differ and a loop when they coincide. Thus the
new graph is oriented.

Consider a new two-step path $v\to u\to x$. It means that
$\rho(v)\to u$ and $\rho(u)\to x$ are original arcs. Applying
incoming containment at $u$ to the first of these arcs yields the
original arc $\rho(v)\to\rho(u)$. Hence
\[
 \rho(v)\to\rho(u)\to x
\]
is an original two-step path. Its intermediate vertex cannot equal
$\rho(v)$, since that would require a loop, and its endpoint cannot
equal $\rho(v)$, since that would require a digon. If $x$ is an exact
second target of $v$ in $D^\rho$, then
$x\notin N_{D^\rho}^+(v)=N_D^+(\rho(v))$. Therefore it is an
original exact second target of $\rho(v)$, proving the inclusion in
\eqref{eq:paper-clone-map-gaps}. The first neighborhood has exactly
the copied original size, so the gap inequality follows.
\end{proof}

\begin{corollary}[Copying a prescribed subset]
\label{cor:paper-subset-copy}
Let $A\subseteq V(D)$, and choose $\rho(a)$ for each $a\in A$ with
$N_D^-(a)\subseteq N_D^-(\rho(a))$. Extend $\rho$ by the identity
outside $A$. Then simultaneous copying gives
\[
 g_{D^\rho}(a)\geq g_D(\rho(a))\quad(a\in A),\qquad
 g_{D^\rho}(v)\geq g_D(v)\quad(v\notin A).
\]
\end{corollary}
\begin{proof}
The extended map satisfies \eqref{eq:paper-clone-containment} at every
vertex, so Theorem~\ref{thm:clone-map} applies.
\end{proof}

The simultaneous formulation is essential for coupled choices of the
map: all containments refer to one graph, and all rows are copied from
that same graph. A sequence that recomputes containment after each
individual copy is a different procedure and is not needed here.

\subsection{A unit-predecessor cover at minimum order}

For the remainder of this subsection, assume that $D$ is all-deficient
and has the smallest possible number of vertices among all
all-deficient oriented graphs. This is a global minimum-order
assumption. For $f\in V(D)$ define
\[
 A_f=\{a\in V(D):a\to f\text{ and }g_D(a)=1\}.
\]
The next obstruction concerns a simultaneous choice of replacement
vertices outside the set being copied.

\begin{lemma}[Unit predecessors cannot all be copied outside their set]
\label{lem:paper-unit-predecessor-obstruction}
For no vertex $f$ is there a map
\[
 \rho:A_f\longrightarrow V(D)\setminus A_f
 \quad\text{such that}\quad
 N_D^-(a)\subseteq N_D^-(\rho(a))\quad(a\in A_f).
\]
In particular $A_f\ne\varnothing$ for every $f$.
\end{lemma}
\begin{proof}
Suppose such a map exists. Extend it by the identity outside $A_f$
and copy the rows simultaneously. The resulting graph $D^\rho$ is
all-deficient by Theorem~\ref{thm:clone-map}.

We claim that every predecessor of $f$ in $D^\rho$ has gap at least
two. An unchanged predecessor was originally a predecessor outside
$A_f$, and hence had positive integral gap at least two; its gap has
not decreased. If a changed vertex $a\in A_f$ still points to $f$,
then its chosen image $b=\rho(a)$ originally pointed to $f$.
Since $b\notin A_f$, we again have $g_D(b)\geq2$, and the copy
inequality gives $g_{D^\rho}(a)\geq2$. There are no other new
predecessors because all rows outside $A_f$ were unchanged.

Delete $f$ from $D^\rho$. Each of its predecessors loses one first
target and can only lose exact second targets among the remaining
vertices; its gap therefore remains at least one. Every other
remaining vertex loses no first target and acquires no exact second
target, so its gap also remains positive. The deleted graph is
nonempty: an all-deficient oriented graph has positive outdegree at
every vertex and therefore at least three vertices. We have produced
a smaller all-deficient graph, contradicting minimum order.

If $A_f$ were empty, the empty map would meet the stated conditions
and the same deletion argument would apply.
\end{proof}

\begin{theorem}[Every vertex receives a maximal pure-unit class]
\label{thm:unit-class-cover}
For every $f\in V(D)$, some maximal incoming class $C$ satisfies
\[
 C\subseteq A_f.
\]
Equivalently, every member of $C$ has original gap one and sends an
original arc to $f$.
\end{theorem}
\begin{proof}
Suppose no maximal incoming class is contained in $A_f$. Each maximal
class then has a representative outside $A_f$. For every $a\in A_f$,
choose a maximal class whose incoming set contains $N_D^-(a)$, and
let $\rho(a)$ be its chosen representative outside $A_f$. Such a
maximal class exists in the finite inclusion poset. The resulting map
satisfies
\[
 N_D^-(a)\subseteq N_D^-(\rho(a)),\qquad
 \rho(a)\notin A_f\qquad(a\in A_f),
\]
contradicting Lemma~\ref{lem:paper-unit-predecessor-obstruction}.
\end{proof}

Call a maximal incoming class \emph{pure-unit} if all its members have
gap one, and let $\mathcal C_1$ be the family of these classes.

\begin{corollary}[A complete-arc cycle of pure-unit classes]
\label{cor:paper-unit-class-cycle}
There are at least three classes in $\mathcal C_1$. Define an oriented
quotient on $\mathcal C_1$ by placing an arc $C\to C'$ when every
vertex of $C$ points to every vertex of $C'$. This quotient has positive
minimum indegree and contains a directed cycle of length at least three.
Selecting one vertex from each class on such a cycle gives an original
directed cycle all of whose vertices have gap one.
\end{corollary}
\begin{proof}
Theorem~\ref{thm:unit-class-cover} makes $\mathcal C_1$ nonempty.
Fix $C'\in\mathcal C_1$ and $f\in C'$. The theorem gives
$C\in\mathcal C_1$ with every member of $C$ pointing to $f$.
Every vertex of $C'$ has the same incoming neighborhood as $f$;
hence all these arcs extend to complete arcs from $C$ to $C'$.
The classes are distinct, since an incoming class is independent.
Thus each quotient vertex has an incoming arc.

The quotient has no loops or digons, by its definition and the
orientedness of $D$. Following predecessors in this finite quotient
produces a directed cycle. Such a cycle has at least three vertices.
Its complete interclass arcs allow any choice of one original vertex
from each cycle class, proving the last assertion.
\end{proof}

Combining this with Corollary~\ref{cor:paper-counterexample-six-incoming-classes},
a minimum-order counterexample would have at least six maximal incoming
classes, at least three of which are pure-unit. Neither assertion says
that every maximal incoming class must be pure-unit.

\subsection{A lexicographic extremal normal form}

Suppose a counterexample exists. Among all all-deficient oriented graphs,
choose $D$ according to the following successive objectives:
\begin{enumerate}
 \item minimize $|V(D)|$;
 \item subject to the first objective, minimize $|A(D)|$;
 \item subject to the first two objectives, maximize
       $\sum_{v\in V(D)}g_D(v)$;
 \item subject to the first three objectives, maximize
       \[
        \Psi(D)=\sum_F
              |\{v\in V(D):N_D^+(v)=F\}|^2,
       \]
       where the sum ranges over the distinct outgoing neighborhoods.
\end{enumerate}
These choices exist conditionally on the existence of a counterexample:
the first minimum is a positive integer, and after fixing it only finitely
many labelled oriented graphs remain. The second objective ranges over
all counterexamples of that order, rather than just subgraphs of one
chosen graph. Call a graph selected in this way \emph{globally extremal}.

\begin{lemma}
\label{lem:paper-extremal-connectivity}
A globally extremal counterexample is strongly connected and minimal
under deletion of any nonempty set of arcs while the vertex set is fixed.
\end{lemma}
\begin{proof}
If it were not strongly connected, choose a sink strongly connected
component. It is a proper nonempty vertex subset with no outgoing arc
to its complement. Both first and second neighborhoods of its vertices
therefore remain unchanged in the induced subgraph. That subgraph
would be all-deficient, contrary to minimum order. If deleting a
nonempty arc set preserved all-deficiency, the resulting graph would
have the same order and fewer arcs, contrary to the second objective.
\end{proof}

\begin{proposition}[Degree and gap monotonicity under incoming containment]
\label{prop:normal-degree}
In a globally extremal counterexample, if $a\ne b$ and
$N_D^-(a)\subseteq N_D^-(b)$, then
\[
 d_D^+(a)\leq d_D^+(b).
\]
If these outdegrees are equal, then $g_D(a)\geq g_D(b)$. In particular,
equal incoming neighborhoods imply equal outdegrees and equal gaps.
If in addition $a\to b$, the missing degrees satisfy
\[
 \mu_D(a)\geq\mu_D(b)+1,
 \qquad \mu_D(v)=|V(D)|-1-d_D^+(v)-d_D^-(v).
\]
\end{proposition}
\begin{proof}
By Lemma~\ref{lem:paper-single-row-copy}, copying the outgoing row
of $b$ to $a$ preserves all-deficiency. It changes the arc count by
$d_D^+(b)-d_D^+(a)$. A negative change would contradict the second
extremal objective, which proves the degree inequality.

If the degrees coincide, the copied graph has the same order and arc
count. Its new gap at $a$ is $g_D(b)$, and all other gaps do not
decrease. Its total gap is thus at least the old total plus
$g_D(b)-g_D(a)$. The third objective excludes a positive increase,
giving $g_D(a)\geq g_D(b)$. When the incoming neighborhoods are
equal, apply both conclusions in both directions.

If $a\to b$, then $a\in N_D^-(b)\setminus N_D^-(a)$, so incoming
containment yields $d_D^-(b)\geq d_D^-(a)+1$. Therefore
\[
 \mu_D(a)-\mu_D(b)
 =d_D^+(b)-d_D^+(a)+d_D^-(b)-d_D^-(a)\geq1.
\]
\end{proof}

\begin{theorem}[Independent incoming classes in an extremal graph]
\label{thm:twin-normal-form}
In a globally extremal counterexample,
\[
 N_D^-(a)=N_D^-(b)\quad\Longrightarrow\quad N_D^+(a)=N_D^+(b).
\]
Consequently every incoming class is an independent module: each outside
vertex has uniform adjacency and orientation to the entire class.
Members of a class have the same first and exact second neighborhoods,
the same gap, and the same missing degree.
\end{theorem}
\begin{proof}
Suppose that $a$ and $b$ have the same incoming neighborhood but
different outgoing neighborhoods $F_a$ and $F_b$.
Proposition~\ref{prop:normal-degree} gives equal outdegrees and equal
gaps. Both single-row copying directions are permitted. They preserve
the first two extremal objectives; their total gaps cannot decrease,
and by the third objective cannot increase. Thus both operations also
preserve the third objective.

Let $r_a$ and $r_b$ be the numbers of vertices whose outgoing
neighborhoods are $F_a$ and $F_b$, respectively. Interchange the names
if necessary so that $r_a\leq r_b$, and copy $b$'s row to $a$.
Every other outgoing row is unchanged. The value of the fourth
objective changes by
\[
 (r_a-1)^2+(r_b+1)^2-r_a^2-r_b^2
       =2(r_b-r_a)+2>0.
\]
This remains the correct formula if the old row class disappears.
The increase contradicts the fourth extremal objective, proving
equality of the outgoing neighborhoods.

An incoming class is independent, as noted above. Equality of both
incoming and outgoing neighborhoods gives uniform adjacency to every
outside vertex, including uniform nonadjacency when no arc is present.
The unions of outgoing neighborhoods over the common first set are
identical for all members of the class. Such a union contains no
member of the class: a path $a\to x\to b$ between two class members
would imply $x\to a$ from their equal incoming sets and hence a
digon. Subtracting the common first set therefore gives identical
exact second sets. Equality of the gaps and missing degrees follows.
\end{proof}

The independent-module conclusion concerns the selected globally
extremal graph. The row-copying inequalities hold without that selection,
but the degree ordering, gap ordering, and final symmetrization use its
successive objectives. In particular, arc-deletion minimality alone
does not supply the tie-breaking properties used in
Proposition~\ref{prop:normal-degree} and
Theorem~\ref{thm:twin-normal-form}.
